\documentclass[11pt]{article}
\usepackage[margin=1in]{geometry}
\usepackage[T1]{fontenc}
\usepackage{lmodern}
\usepackage{microtype}
\usepackage{booktabs}
\usepackage{enumitem}
\usepackage{amsmath}
\usepackage{xcolor}
\usepackage[hidelinks]{hyperref}

\definecolor{eyeblue}{HTML}{203B78}
\definecolor{eyegreen}{HTML}{1B756B}
\setlist{nosep,leftmargin=*}
\setlength{\parskip}{0.55em}
\setlength{\parindent}{0pt}

\title{\textcolor{eyeblue}{EyeExplore User Study}\\
\large Final Hypotheses and Evaluation Plan}
\author{}
\date{}

\begin{document}
\maketitle

\section{Key terms}
This report proposes two testable ideas about using different views. The first asks
whether looking at two related views together helps people answer analysis
questions. The second asks whether this approach helps novice users perform
closer to experienced users. 

\begin{itemize}
  \item A \textbf{view} is one display of gaze data, such as an attention map,
  scanpath, or timeline.
  \item In the \textbf{one-view layout}, a participant sees one display at a
  time and may switch to another display when needed.
  \item In the \textbf{two-view layout}, two displays of the same study data
  are visible side by side. The views are chosen because they show related
  information.
  \item A \textbf{novice} is a person with little or no previous experience
  analyzing eye-tracking data. An \textbf{expert} is a person with regular,
  substantial experience doing this analysis. The study team must set the exact
  experience rule before recruiting.
\end{itemize}

\section{The two final hypotheses}

\subsection*{H1 - Seeing two related views helps people analyze gaze data}
\textbf{Hypothesis:} People who can see two related views at the same time will
answer gaze-analysis questions more accurately than people who can see only
one view at a time. They may also finish the questions faster.

\textbf{What this means in practice:} Suppose a researcher wants to know both
where people looked and the order in which they looked there. An attention map
can show the main areas of interest, while a scanpath or AOI timeline can show
the order and timing of gaze. With one view, the researcher must inspect one
display, switch views, and remember what was shown. With two views, the
researcher can compare the related evidence side by side. If the two displays
make the same pattern easier to check, the researcher should make fewer wrong
answers.

The same idea applies to video. A video shows the scene and the viewer's gaze
movement. A synchronized K-Coefficient timeline shows the attention value at
the corresponding time. Seeing the video and timeline together may help a
person connect a change in the measured attention value to what was happening
in the video. This is one example of two related views; the study is not
claiming to test all 13 visualization types separately.

\textbf{What result would support H1:} Within the novice group, the two-view
group should get more answers right than the one-view group in less time. The same
comparison is made within the expert group. Faster completion in the two-view
group would add support, but accuracy can be the main result. If the groups score
about the same, or the one-view group scores higher, the study does not support
the prediction of proposed hypothesis.

\subsection*{H2 - Two-view EyeExplore can bring novice results closer to expert results}
\textbf{Hypothesis:} After the same short introduction, novice users working
with two related views will answer the shared analysis questions nearly as
accurately as expert users working with the same two-view setup.

\textbf{What this means in practice:} Experts have more experience interpreting
gaze patterns. The question is whether EyeExplore's two-view layout and a short
introduction help novices make correct findings that are close to the findings
made by experts. This does not assume that the two groups will have identical
speed or experience. It predicts that their accuracy on the same questions can
be close.

For this study, ``nearly as accurately'' means that the novice group's score on
the shared questions is no more than some X percentage points below the expert
group's score (X can be dependent on dataset or level of question asked , so we need to work on X). For example, if experts answer 80\% correctly, novices would
need to answer at least 80-X\% correctly to meet this proposed rule. The X-point
margin is a study decision, not a known fact; the team should approve it before
collecting data. The results must also be precise enough to rule out a larger
gap. If the sample is too small to make that conclusion, report the result as
unclear rather than saying the groups are equal.

\section{Study design}
Use a 2 by 2 design. The two factors are user experience (novice or expert) and
screen layout (one view or two views side by side).

\begin{center}
\begin{tabular}{lll}
\toprule
Group & Experience & Layout \\
\midrule
Set 1 & Novice & One view at a time \\
Set 2 & Novice & Two related views at the same time \\
Set 3 & Expert & One view at a time \\
Set 4 & Expert & Two related views at the same time \\
\bottomrule
\end{tabular}
\end{center}

Each participant completes one set only. Within novice volunteers, assign
participants at random to Set 1 or Set 2. Do the same within expert
volunteers. This avoids participants seeing the same questions twice and
remembering the answers. Use the same study files, stimuli, records, and
question wording in the two layout groups within each experience level.

\textbf{Important interface requirement.} The study needs a real one-view
condition: one visualization occupies the work area, and the second view is
not visible. Two identical views shown in two panels do not count as one view.
If the current EyeExplore app has no one-view mode, add or configure a controlled
study mode before recruitment. Do not manually cover a panel differently for
different participants.

\section{How the questions evaluate the hypotheses}
Every booklet has four scored multiple-choice questions:

\begin{itemize}
  \item Questions 1 and 2 test H1. The novice questions use simpler gaze
  comparisons. The expert questions use more complex comparisons. Within each
  experience level, the one-view and two-view groups receive the same tasks.
  Thus, compare Set 1 with Set 2 and Set 3 with Set 4 for H1.
  \item Questions 3 and 4 are identical in all four sets. They are the shared
  anchor questions for H2. Compare novice and expert scores on these questions,
  especially Set 2 with Set 4 for the stated experience-gap claim.
\end{itemize}

All questions have one correct answer. Options include an equality or
``neither/both'' choice where appropriate. The answer key is based on fixed
data and explicit rules, not on a researcher's opinion about what a plot looks
like.

\section{Procedure}
\begin{enumerate}
  \item Recruit novice and expert users using definitions written before the
  study. For example, novices have not analyzed eye-tracking data before;
  experts have at least two years of regular eye-tracking analysis experience.
  Record experience separately from questionnaire answers.
  \item Give every participant the same five-minute introduction. Explain how
  to select the dataset, stimulus, records, and visualization. Do not teach the
  answers to the study tasks.
  \item Give each participant the booklet set assigned to their group. Read
  instructions as written and do not hint at a correct answer. Allow about 25
  minutes for the four tasks in the set.
  \item Record correctness and task completion time for each question. Use the
  same start and stop rule for everyone (start when the task screen is ready;
  stop when the participant submits an answer).
  \item Record technical failures separately. If a required file or view does
  not load, pause that participant's task and do not count it as an incorrect
  answer.
\end{enumerate}

\section{How to score and compare results}
Score a correct answer as 1 and an incorrect answer as 0. A question left blank
after the set time limit is incorrect; a documented technical failure is
missing data, not an incorrect answer. The primary outcome is the percent of
questions answered correctly. Also report median time per question and the
number of technical failures.

For H1, compare accuracy and time between one-view and two-view groups
separately for novices and experts. Then report the overall layout effect. In
simple terms, calculate how often each group answered correctly and how long
they took. If the study has enough participants, a statistician can also use a
model that checks whether the layout effect changes between novice and expert
users. A higher score in the two-view condition supports H1. Faster time is
supporting evidence, not a requirement for H1 if accuracy improves.

For H2, compare the two shared questions between novice and expert users in the
two-view condition. Report the score difference in percentage points and show a
95\% confidence interval. This interval is a range around the measured
difference that shows how uncertain the estimate is. H2 is supported only if
the data are precise enough to show that the novice score is within the
pre-agreed 10-point margin. If the range is too wide, call the result
inconclusive rather than claiming that the groups are equal.

\textbf{Sample size matters.} Four questions per person give a short pilot
measure, not a strong final measure of expertise. A publication-level study
should include more matched items or repeated, equivalent task cases and use a
sample-size calculation before recruitment. Keep the four-question booklet as
the starter version, and treat its findings as preliminary unless the study is
expanded.

\end{document}
